\documentclass[11pt]{article}
\usepackage{amsmath,amssymb}
\title{A quantitative harmonic approximation for original long heat legs}\author{}\date{}
\begin{document}\maketitle
\paragraph{Input and scope.}
Use the original gaussian coarse kinetic interface calculation and the harmonic remainder and heat equation contacts calculation. In the longest time completed contact packet calculation, the
longest-time packet differentiates $C(G-G_s,g-g_s)$; its difference
from the shortest packet is the stated total derivative with zero
endpoints. The ghost difference has the correct signs, and the
Euler weights are $C_1+2C_2+3C_3$ for the uncut insertion and
$-C_2-2C_3$ for the remainder. Thus the isolated late Haar tensor of the harmonic remainder and heat equation contacts calculation
is combined with its ultraviolet partners in this rearrangement.
This does not prove a form bound for the resulting full tensor.

The following new estimate concerns the ORIGINAL flat relative
one-form complex, before any tree-coordinate norm substitution.
It quantifies an actual long-leg approximation; no nonlinear measure
or generic effective density is substituted. Color copies are direct
sums and change constants only by their literal multiplicity.

\paragraph{1. Complete subblock residual with the Gauss constraint.}
Let $H=d^*d$ on the global cube of side $L$, with relative tangential
boundary conditions. Let $P$ be its physical projection, and take a
family of complete subblocks of side $\ell$ with disjoint free
interiors $J_B$ and retained complete boundary skeleton. All original
plaquettes incident to an interior link belong to its block.
For any global field $w$, define
\[
 q_B=(Hw)|_{J_B},\qquad r_B=H_B^+q_B,\qquad h_B=w|_B-r_B.       \tag{1}
\]
Here $r_B$ has zero boundary and $H_B$ is the complete original
relative Maxwell operator of that block. Crucially $q_B$ is in its
physical range: an interior gauge gradient $D_B\xi$ extends by zero
to a global gauge gradient, so
\[
 \langle D_B\xi,q_B\rangle=\langle D\widetilde\xi,Hw\rangle=0.
\]
Consequently $H_Br_B=q_B$, and $h_B$ is the true local harmonic
extension of the retained boundary, modulo harmless pure gauge.
Its energy is orthogonal to that of $r_B$. No gauge mode is inverted.

On the original relative product complex the completed Laplacian
$\Delta_{1,B}=H_B+D_BD_B^*$ has one Neumann and three Dirichlet
coordinate factors in each edge component. Its smallest eigenvalue
is at least $12\sin^2(\pi/(2\ell))\ge12/\ell^2$.
On co-closed vectors it equals $H_B$. Hence
\[
 \boxed{\sum_B\|dr_B\|^2
 =\sum_B\langle q_B,H_B^+q_B\rangle
 \le\frac{\ell^2}{12}\sum_B\|q_B\|^2
 \le\frac{\ell^2}{12}\|Hw\|^2.}                              \tag{2}
\]
The last inequality uses disjoint interiors, not independent random
blocks. All norms are literal full-link angular norms. Tree coordinates
can be used to compute the same quotient, but their Euclidean norm
cannot replace these norms silently.

\paragraph{2. Actual long covariance and the gained time factor.}
Let $W_s$ be the centered physical Gaussian with covariance
$T_s=H^+e^{-sH}P$. This is the original infrared covariance appearing
in the completed transverse packet, not a new Wilson law.
Since $HT_sH=He^{-sH}P$, (2) proves
\[
 E\sum_B\|dr_B(W_s)\|^2
 \le\frac{\ell^2}{12}\operatorname{Tr}_P(He^{-sH}).             \tag{3}
\]
The product spectrum and the elementary eigenvalue estimate
$\lambda e^{-s\lambda}\le C(1+s)^{-1}e^{-s\lambda/2}$
(the lattice spectrum is bounded) give the uniform estimate
\[
 \boxed{E\sum_B\|dr_B(W_s)\|^2
 \le C\ell^2\,L^4(1+s)^{-3}e^{-cs/L^2}.}                     \tag{4}
\]
The same constants apply to rectangular boxes of bounded aspect ratio,
with the evident volume and longest-side replacement. The ordinary
long-leg energy count is bounded by
$CL^4(1+s)^{-2}e^{-cs/L^2}$. Thus (4) contains one additional
$\ell^2/(1+s)$ factor in the upper bound. This is not a claim that the
ratio of the exact two traces is bounded by $C/(1+s)$ for arbitrarily
late times; their exponential rates can be adjusted as displayed.
For meaningful approximation use blocks shorter than the diffusion
length. No physical-volume root-count multiplication was used.

\paragraph{3. What this does and does not buy for the completed insertion.}
The global tensor in the longest time completed contact packet calculation remains the SAME global uncut tensor;
it is not a collection of independent block tensors. For any linear
assembly $\mathcal H w$ of the actual local harmonic pieces and the
retained skeleton, write $w=\mathcal H w+r$, with $r$ supported in the
disjoint interiors. Then the exact quadratic insertion is
\[
 w^*Kw=(\mathcal H w)^*K\mathcal H w
             +2(\mathcal H w)^*Kr+r^*Kr.                      \tag{5}
\]
The assembled $r$ need not be globally co-closed: local normal currents
can end on retained skeleton vertices. Formula(5) is a full-link
identity. For a physically projected insertion use $Pr$ and
$P\mathcal H w$ instead; their sum is $w$ for physical $w$, and their
curl energies are unchanged, but the global projection is nonlocal.
It cannot be treated as a bounded local connector.

Equation(4) controls the classical energy of $r$, not $K$ acting on it.
The harmonic piece retains ALL boundary modes, not only the 18
constant-curvature modes. The bound of the complete harmonic kinetic matching calculation and the complete constant curvature coarse sector calculation cannot be applied to it
without a further estimate. A self-energy form bound or a completed
signed current identity is still required for the cross and residual
terms in (5). Nonlocal cross-block entries of $K$, ghost marking and
the multiple-long-line remainder of the longest time completed contact packet calculation stay present.

The underlying exact Gaussian skeleton transition is the Schur law of the original gaussian coarse kinetic interface calculation,
with its raw determinant and inherited global/exterior data. Estimate
(2) uses the same original local Maxwell problem; it neither changes
that law nor removes its scalar. Background source derivatives of
$P,H$, the harmonic map and its boundary data remain separate partners.
No bound on those differentiated versions is inferred from (4).

\paragraph{4. Verification.}
The finite calculation uses a literal global $L=4$ original relative complex and
its 16 side-2 blocks. On full (not tree-deleted) link spaces, exact
rational arithmetic checks $D_B^*q_B=0$, $H_Br_B=q_B$, $D_B^*r_B=0$,
and (2) for a nontrivial rational field. Local inversion uses
$H_B+D_BD_B^*$ and the range identity, so no singular gauge inverse is
hidden. This checks the incidence/interface, while the product-spectrum
argument proves the stated size dependence. It is not a check of a
nonlinear two-loop remainder or S94 closure.
\end{document}
